
\subsection{Termination, Soundness and Completeness}\label{sec:proofs} 
We now establish termination, soundness, and completeness of the tableau procedure.
Termination follows from finiteness: the input formula is finite and zones are compared modulo max-constant abstraction, so only finitely many distinct nodes arise. Soundness is shown by turning a {successful marking} of the AND--OR graph into a witnessing TTS -- the branching structure of the marking supplies the discrete successors demanded by the existential markers, while the context $\mathcal{X}$ supplies the universal ones -- and then unfolding it into a non-Zeno run. Completeness follows by guiding the construction with a satisfying model so as to produce a successful marking. We present the results here, with detailed proofs in Appendix~\ref{sec:appendix}.

The key to establish termination is the finiteness of distinct tableau nodes, which follows from the finiteness of the input formula and the use of abstraction to ensure a finite search space. We formalize this in the following lemma.

\begin{lemma}[Finiteness of nodes]\label{lem:finiteness}
%Fix an input formula $\varphi$ of length $n = |\varphi|$ and let $M$ be the maximal integer constant
%occurring in any timing interval of $\varphi$.
%Assume zones are maintained in canonical form modulo $\cong_M$. 
%Then only finitely many distinct tableau nodes can arise during the
%construction and the number of nodes is bounded by $2^{O(n \log n)} \cdot M^{O(n)}$

Let $\varphi$ be an input formula of size $n=|\varphi|$ and $M$ the largest constant in its timing intervals. With zones normalised modulo $\cong_M$ and one formula-clock per timed-until subformula, the number of clocks is $K=O(n)$, and the number of distinct tableau nodes is finite, bounded by $2^{O(n\log n)}\cdot M^{O(n^2)}$. Measuring timing constants in unary (the standard convention for dense-time logics), this bound is singly exponential in $|\varphi|$.

\end{lemma}

\begin{proof}
    In Appendix~\ref{sec:appendix}.
\end{proof}

From the finiteness of nodes, we can establish termination of the procedure.

\begin{theorem}[Termination]\label{thm:termination}
For every input formula $\varphi$, the depth-first tableau procedure terminates.
Equivalently, every tableau branch is finite.
\end{theorem}

\begin{proof}[Sketch]
By Lemma~\ref{lem:finiteness}, only finitely many distinct nodes can occur, so any infinite branch would contain a repeated node; upon such a repetition the loop rule either accepts (back-edge) or closes, so no branch is unfolded past a repeat and every branch is finite. Each node has finite out-degree (at most five children from an expansion rule, at most $|\Gamma|$ from successor generation), so by König's lemma the entire memoised construction is finite.
\end{proof}
\noindent
Having established termination, we can now establish an upper bound on the complexity of the procedure.


\begin{theorem}[EXPTIME-completeness]\label{thm:complexity}
Satisfiability of the non-punctual fragment decided by the procedure is EXPTIME-complete; with timing constants in unary the procedure decides it in deterministic exponential time in $|\varphi|$.
\end{theorem}
\begin{proof}[Sketch]
{Membership.} Lemma~\ref{lem:finiteness} bounds the number of {distinct} nodes. The bound is realised only by {memoising} visited nodes, so that each distinct node is expanded once and the construction materialises the finite AND--OR {graph} rather than its unfolding; the naive ancestor-only depth-first search may otherwise explore a tree that is doubly exponential. A successful marking must make {every} infinite path good (fulfilling and time-progressing); since markings retain {all} AND-children, this is a {universal} path condition over the AND-branching, not a per-SCC or per-cycle property -- an SCC that AND-forks into one good and one bad cycle must be rejected. Existence of a successful marking is therefore the winning condition of a two-player game on the finite AND--OR graph: the proponent resolves the OR-choices, the opponent resolves the AND-branches and the path, and the proponent wins a play iff it is an accepting lasso (every eventuality fulfilled, and time-progressing by Lemma~\ref{lem:nonzeno}). This is the standard branching-time acceptance -- a Streett condition with $k=O(n)$ pairs (the eventualities and the delay-bounding clocks) -- solved by the usual nested-fixpoint elimination over the AND--OR graph~\cite{Clarke99,tableau_ctl}. Solving a $k$-pair Streett game over $m$ nodes takes $m^{O(k)}$ time; here $m=2^{O(n\log n)}\cdot M^{O(n^2)}$ is already singly exponential in $|\varphi|$ (with constants in unary), so $\log m=\mathrm{poly}(|\varphi|)$ and $m^{O(k)}=2^{O(k\log m)}=2^{\mathrm{poly}(|\varphi|)}$ remains singly exponential. Hence membership is in deterministic exponential time.
{Hardness.} Restricting all intervals to $[0,\infty)$ the procedure decides CTL satisfiability, which is EXPTIME-hard~\cite{Clarke99}; hence EXPTIME-completeness.
\end{proof}

\begin{lemma}[Non-Zeno characterisation]\label{lem:nonzeno}
Let $N_i\equiv N_j$ be a candidate loop. Its $\omega$-unfolding admits a non-Zeno (time-divergent) run iff the loop is {time-progressing}: it admits positive total elapsed time {and} every clock not reset on the segment is extrapolated above $M$ on it. When it is, a non-Zeno run exists in which every iteration elapses the same positive amount.
\end{lemma}
\begin{proof}[Sketch]
($\Rightarrow$) If some clock $x$ is not reset on the segment yet keeps a finite upper bound $c\le M$, then along the loop $x$ only increases while staying $\le c$, so the accumulated delay is at most $c-x_{\text{entry}}<\infty$ and every unfolding is Zeno; likewise if the loop can only elapse zero time. ($\Leftarrow$) If the loop is time-progressing, fix a realisable $\delta>0$ for one traversal. Since $N_i\equiv N_j$ the segment repeats structurally, and every delay-bounding clock is either reset within the loop (its budget replenished) or saturated above $M$ (no bound); hence $\delta$ is realisable every iteration and $\sum_i t_i=\infty$. Full proof in Appendix~\ref{sec:appendix}.
\end{proof}



\begin{remark}
    A one-pass tableau for TCTL could be derived by extending the context-injection methodology of for e.g. \cite{tableau_ctl}. However, this approach incurs a double-exponential blowup in the size of the tableau due to the combinatorial explosion of contexts. Furthermore, accurately capturing non-Zeno requirements within a purely syntactic context is non-trivial. Cycle evaluation is an efficient way to enforce non-Zenoness; thus, by tracking eventualities semantically alongside time divergence, we bypass the combinatorial explosion of context-injection and maintain a tractable state space.

\end{remark}


\begin{remark} The worst-case $M^{O(n^2)}$ zone bound is triggered specifically by adversarial, multi-clock specifications—for instance, when a global eventuality tracks $M$ while a local invariant forces high-frequency resets of a secondary clock. In such cases, maintaining canonical zones forces the tableau to fragment the continuous timeline to track precise clock differences. However, for many standard time-bounded properties, such as $E(\top U_{[0, M]} p)$, the procedure allocates a single clock and evaluates the entire interval efficiently via a single symbolic zone intersection. Because zones serve as symbolic representations of infinite sets of clock valuations, the tableau naturally absorbs large values of $M$ in practice without explicit enumeration, making the algorithm highly efficient for typical verification workloads.

\end{remark}




To establish soundness we first analyse the shape of a successful marking. Following standard conventions in finite-state verification~\cite{Clarke99}, the infinite paths of the marking are lasso-shaped: a lasso is a finite prefix followed by an infinitely repeating finite cycle, $\pi=\pi_{\mathsf{pref}}\cdot \pi_{\mathsf{loop}}^\omega$. The following lemma guarantees that {every} infinite path of a successful marking is such a lasso whose loop segment satisfies the acceptance condition; this is what lets us realise the universal obligations along all branches, not just one.





\begin{lemma}[Lasso extraction]\label{lem:lasso}
Let $\mathcal{M}$ be a successful marking. Then $\mathcal{M}$ is finite, and every infinite path of $\mathcal{M}$ is a lasso $\pi=\pi_{\mathsf{pref}}\cdot \pi_{\mathsf{loop}}^\omega$ whose loop segment is accepting: every postponed eventuality is fulfilled on $\pi_{\mathsf{loop}}$ and $\pi_{\mathsf{loop}}$ is time-progressing.
\end{lemma}

\begin{proof}
    In Appendix~\ref{sec:appendix}.
\end{proof}

A successful marking is thus a finite symbolic representation of a (branching) family of infinite timed executions, each satisfying all postponed eventualities. From the marking we extract a TTS that witnesses satisfiability, which we use to establish soundness.

The truth induction in soundness proceeds over the {discrete} nodes, yet the semantics (Definition~\ref{def:sat}) quantifies over {all} positions, including those strictly inside a delay. The following lemma bridges the gap; it is where the integer-hyperplane split of successor generation earns its keep.

\begin{lemma}[Region stability]\label{lem:stability}
Along a delay step $(s,\nu)\xrightarrow{t}(s,\nu+t)$ that crosses no integer hyperplane $x=b$ ($b\le M$) of any clock, the truth of every formula in $Sub(\varphi)$ is constant over the {relatively-open} segment $(s,\nu+t')$, $t'\in(0,t)$; the segment endpoints, where a clock may sit on a hyperplane, are themselves nodes (materialised by the integer-split) and are handled by the strictness flags of the interval bounds.
\end{lemma}
\begin{proof}[Sketch]
A delay segment that crosses no integer hyperplane $x=b$ ($b\le M$) lies entirely within a single clock {region}: no clock changes integer part below $M$ and no fractional part wraps, so all configurations $(s,\nu+t')$, $t'\in[0,t]$, are region-equivalent. TCTL truth is constant on a region (the standard region theorem~\cite{AlurCourcoubetisDill1990,Bouyer05TimedAutomata}); hence every formula in $Sub(\varphi)$ has constant truth across the segment. The boundary-split (Section~\ref{sec:tableau}) makes the segment endpoints nodes, so establishing truth there establishes it at every interior position. (Diagonal-freeness and integer endpoints suffice; non-punctuality is not needed for this lemma.)
\end{proof}

The soundness argument for the timed modalities rests on identifying the {computed} satisfaction set $\mathrm{Sat}(\cdot)$ with its {semantic} counterpart in the model that the marking builds. This is the load-bearing step, so we state it explicitly.

\begin{lemma}[Correctness of $\mathrm{Sat}$]\label{lem:sat}
Let $\mathcal{M}$ be a successful marking and $\mathcal{T}$ the model read off its pursuit branches (Section~\ref{sec:proofs}, soundness). For every timed modality $\xi$ demanded at a node $s$ and every $\nu\in Z_s$, $\;\nu\in\mathrm{Sat}(\xi)\iff \mathcal{T},(s,\nu)\models\xi$.
\end{lemma}
\begin{proof}[Sketch]
The runs of $\mathcal{T}$ from $(s,\nu)$ are exactly the pursuit paths of $\mathcal{M}$ from the check rooted at $s$: the discrete edges of $\mathcal{T}$ are the pursuit's $\EX$-/$\AX$-successors, and the left-drop continuations contribute no run. The check computes $\mathrm{Sat}(\xi)$ as the least- (until) or greatest- (release) fixpoint of the one-step unfolding~\eqref{tag:teu}/\eqref{tag:ter}: the stop/in-region branch contributes valuations with an immediate $\psi$-witness at $x_\xi\in\kI$, the go/before branch the time-predecessors of the successors' set. This is precisely the fixpoint characterisation of the matching clause of Definition~\ref{def:sat} over those runs (region stability, Lemma~\ref{lem:stability}, covers the dense positions). For the universal variants the unfolding quantifies over {all} successors and the deadline/entry invariant -- placed on the pursuit edges, which are the {only} edges of $\mathcal{T}$ -- forces every run through the window, so the fixpoint is the genuinely universal set. Hence $\mathrm{Sat}(\xi)=\{\nu : \mathcal{T},(s,\nu)\models\xi\}$.
\end{proof}

\begin{theorem}[Soundness]\label{thm:soundness}
If the tableau succeeds on $\varphi$, then $\varphi$ is satisfiable.
\end{theorem}

\begin{proof}[Sketch]
From a successful marking $\mathcal{M}$ (Lemma~\ref{lem:lasso}) we build a finite TTS whose locations are the nodes of $\mathcal{M}$ and whose discrete edges are the AND-children retained by $\mathcal{M}$: every existential marker $\EX\psi$ thereby obtains a successor satisfying $\psi$, and every node's context $\mathcal{X}$ is imposed on {all} its successors, so the universal obligations ($\AX$, $\mathbf{A}$-until, $\mathbf{A}$-release) hold across the branching rather than vacuously along a single path. Zone non-emptiness along $\mathcal{M}$ yields, by a feasibility argument (Lemma~\ref{lem:feasible}, appendix), a coherent assignment of delays and valuations threading the time-elapse and reset steps; the time-progress condition on every loop (Lemma~\ref{lem:nonzeno}) lets us schedule strictly positive per-cycle delays, so all runs of the model are time-divergent. A structural induction on $\varphi$ -- the accepting-loop condition discharging the least-fixpoint (until) eventualities and co-induction discharging the greatest-fixpoint (release) obligations -- shows the model satisfies $\varphi$ at $(s_0,\nu_0)$.
\end{proof}

%We now prove completeness, showing that if $\varphi$ is satisfiable, then the tableau procedure successfully finds a proof. The key idea is to construct a branch guided by the satisfying model, ensuring that the tableau expansion and successor generation steps are consistent with the model's behavior. Because the model satisfies $\varphi$, it must eventually fulfill all eventualities, which guarantees that the constructed branch will be successful.

It remains to show that the tableau procedure is complete. Intuitively, the tableau construction is expressive enough to reconstruct any satisfying execution. Thus, whenever a formula admits a model, the tableau should be able to realize a corresponding successful branch. We now state the completeness result formally.


%Suppose $\varphi$ is satisfiable in some TTS, then we can construct a tableau branch guided by this model. At each expansion step, we choose the successor node that corresponds to the behavior of the model. 


\begin{theorem}[Completeness]\label{thm:completeness}
If $\varphi$ is satisfiable, then the tableau succeeds on $\varphi$.
\end{theorem}


\begin{proof}[Sketch]
Let $\mathcal{T},(s_0,\nu_0)\models\varphi$ with $\mathcal{T}$ non-Zeno. We build a marking by following $\mathcal{T}$'s branching: at each OR-node we select the child consistent with $\mathcal{T}$; at each successor-generation (AND) node we realise {every} demanded existential successor by a corresponding successor of $\mathcal{T}$ and recurse on all of them, while the universal context $\mathcal{X}$ holds at each because $\mathcal{T}$ satisfies the $\AX$-obligations. Every chosen zone intersection is non-empty, witnessed by the model's valuations. A timed until demanded by a persistent context is verified from each state of $\mathcal{T}$ on its own dedicated-check branch (Section~\ref{sec:tableau}), with its formula-clock reset there, so staggered demands are witnessed {independently} -- following $\mathcal{T}$'s own (possibly branching) witnesses -- and the dense family collapses to finitely many region-classes. As only finitely many distinct nodes exist, every path of the marking eventually repeats; because $\mathcal{T}$ satisfies each until within finite time, every postponed eventuality is fulfilled on the cycle, and because $\mathcal{T}$ is non-Zeno every cycle is time-progressing (Lemma~\ref{lem:nonzeno}). Hence the marking is successful, and completeness holds for {all} non-punctual TCTL.
\end{proof}



\begin{remark}
  We restrict the evaluation of temporal modalities exclusively to non-Zeno paths, which is a standard assumption in the literature on real-time logics and timed transition systems~\cite{AlurDill1994,BouyerHR2008}. This restriction is strictly motivated by the fact that Zeno behaviors are not physically realisable. Permitting such time-convergent executions can lead to spurious witnesses for satisfiability, as it allows a model to vacuously fulfill time-bounded safety or release obligations simply by stalling the progression of physical time. However, the framework supports {both} semantics with a single switch: dropping acceptance condition~(ii) (the time-progress test) while keeping~(i) yields a sound and complete procedure for the unrestricted semantics over all mathematically valid paths, including Zeno executions. Soundness and completeness for that variant follow from the same proofs with the non-Zeno scheduling step (Lemma~\ref{lem:nonzeno}) removed; the non-Zeno restriction is thus an optional, well-motivated refinement rather than a load-bearing assumption.
\end{remark}







%

%
%
%During depth-first exploration, whenever a node matches an ancestor on the
%active branch, a candidate loop is detected. Let $N_i$ and $N_j$ with $i<j$
%such that $N_i \equiv N_j$. The loop from $N_i$ to $N_j$ is {accepting}
%iff every eventuality $e \in \eventualities_i$ is fulfilled on the segment
%$[i,j)$, that is, there exists $k$ with $i \le k < j$ such that the stop
%decomposition of $e$ is applied and, in the timed case, the corresponding
%interval constraint is enforced by a nonempty zone refinement.
%
%If the loop is accepting, the branch is declared successful.
%Otherwise it is rejecting and closes at $N_j$.
%
%
%The acceptance condition enforces precisely the fairness requirement for
%least-fixpoint operators. If an eventuality postponed at the entry of the
%cycle were never fulfilled along the cycle, its infinite repetition would
%yield an infinite path violating the semantics of until. Conversely, if every
%postponed eventuality is fulfilled within the cycle, then each iteration of
%the cycle discharges all active obligations, and its infinite unfolding
%satisfies the formula.
%
%Termination follows from finiteness of abstractions.
%The sets $\Gamma$ and $\mathcal{X}$ range over the finite closure of the
%input formula. The set of clocks is bounded by the number of timed
%eventualities. Finally, zones are compared modulo semantic equality and are
%restricted by the maximal constants occurring in the formula; hence only
%finitely many canonical zone representatives exist. Therefore only finitely
%many distinct abstractions $\mathsf{abs}(N)$ can arise along a branch,
%and the loop rule guarantees termination.
%
%
%A successful open branch induces a finite symbolic lasso, consisting of a
%finite prefix followed by an accepting cycle. This lasso can be unfolded
%into a witness timed transition system: tableau nodes correspond to abstract
%states, discrete successor rules induce discrete transitions, and time-elapse
%steps correspond to delay transitions justified by zone closure. By
%construction of the expansion rules and the loop acceptance condition,
%the extracted model satisfies the initial formula.
%Soundness and completeness follow by induction on formula structure,
%together with correctness of zone operations and the fairness argument for
%ultimately periodic models.




